% DENSITYPOSFAIL1
\documentclass[11pt]{article}
\usepackage[margin=.9in]{geometry}
\usepackage{amsmath,amssymb,amsthm}
\newtheorem{proposition}{Proposition}
\newtheorem{lemma}{Lemma}
\begin{document}
\begin{center}
{\Large DENSITYPOSFAIL1: Barriers for positive test families on the density-one slice}\\[4pt]
\textbf{Date: 2026-10-04}\\
\textbf{Status: CANDIDATE pending independent review}\\
\textbf{P versus NP status: OPEN}
\end{center}

\section{Setup}
Let \(M=8\binom q3\) and let the primary exact layer be \(|A|=q\), as in SLICEDENSITY1. Uniform exact-\(q\) formulas are SAT with high probability, while exact-layer SAT remains NP-complete.

A distributional approximation argument now needs a positive distribution supported on exact-\(q\) UNSAT formulas that cannot be separated from the uniform SAT layer by a small monotone circuit.

\section{General recognizability barrier}

\begin{proposition}[P/POLY-POSITIVE-BARRIER]
Let \(\mathcal Y_q\subseteq\binom{[M]}q\) be any positive test family. Suppose membership in \(\mathcal Y_q\) has polynomial-size unrestricted Boolean circuits. Then there is a polynomial-size monotone circuit \(G_q\) such that on the exact layer
\[
G_q(A)=1\iff A\in\mathcal Y_q.
\]
In particular, \(G_q\) perfectly separates \(\mathcal Y_q\) from every disjoint negative test family on the same layer.
\end{proposition}

\begin{proof}
Define the total slice extension
\[
H_q(A)=
\begin{cases}
0,&|A|<q,\\
1[A\in\mathcal Y_q],&|A|=q,\\
1,&|A|>q.
\end{cases}
\]
Equal-weight inputs are pairwise incomparable, so \(H_q\) is monotone. A polynomial unrestricted circuit computes the Hamming weight, tests equality/inequality with \(q\), and on the equality layer runs the membership circuit for \(\mathcal Y_q\). Valiant's slice theorem converts this into a polynomial-size monotone circuit computing \(H_q\).
\end{proof}

\paragraph{Consequence.}
A positive family chosen because it has a simple polynomial-time recognition theorem cannot support a distributional lower bound against all polynomial-size monotone circuits. This directly warns against fixed-deficiency minimally unsatisfiable families: fixed-deficiency minimal-unsatisfiability recognition is known to be polynomial/FPT for fixed deficiency.

\section{Hostile test: hidden density-six active core}

Consider the tempting planted distribution:
\begin{enumerate}
\item choose a uniform active set \(S\subseteq[q]\) of size \(q/6\);
\item choose exactly \(q\) proper 3-clauses uniformly from clauses supported on \(S\);
\item condition on UNSAT.
\end{enumerate}
The conditioning changes probability by only \(\exp(-\Omega(q))\), because the active formula has clause density six and
\[
\Pr[\mathrm{SAT}]\le 2^{q/6}(7/8)^q.
\]

Let \(d_i(A)\) be the number of selected clauses containing variable \(i\), and define
\[
H_6(A)=|\{i:d_i(A)\ge6\}|.
\]

For the uniform density-one layer, the degree of a fixed variable converges to a Poisson random variable of mean three, so
\[
\frac{H_6}{q}\to
p_-:=\Pr[\mathrm{Pois}(3)\ge6]
\approx0.0839179420.
\]
For the planted active-core distribution, active-variable degree converges to \(\mathrm{Pois}(18)\), while inactive variables have degree zero. Hence
\[
\frac{H_6}{q}\to
p_+:=\frac16\Pr[\mathrm{Pois}(18)\ge6]
\approx0.1666126678.
\]
Sampling without replacement has bounded differences, so both statistics concentrate exponentially around these distinct constants. Therefore any fixed threshold between them, for example
\[
H_6\ge0.12q,
\]
distinguishes the planted positive distribution from the uniform density-one distribution with error \(\exp(-\Omega(q))\).

\begin{lemma}[Polynomial monotone degree-profile separator]
The predicate \(H_6(A)\ge0.12q\) has polynomial-size monotone circuits.
\end{lemma}
\begin{proof}
For each variable \(i\), the predicate \(d_i\ge6\) is the OR over all six-element subsets of clause coordinates incident to \(i\), with an AND for each subset. There are \(O(q^2)\) incident coordinates, so this has polynomial size because six is constant. The resulting \(q\) indicator bits can be thresholded by a polynomial-size monotone sorting network, whose comparators use AND/OR.
\end{proof}

Thus the naive hidden high-density core is \textbf{FALSIFIED/RETIRED} as a positive test distribution.

\section{Design requirements for DENSITY1-POSITIVE1}

A serious positive distribution should satisfy all of the following:
\begin{enumerate}
\item exact clause count \(q\) and UNSAT support;
\item no known polynomial-size membership/recognition circuit for its support;
\item local variable-degree histogram asymptotically matching the uniform exact-\(q\) layer;
\item matching low-order clause-overlap and codegree statistics, not only one-coordinate marginals;
\item no constant-size or fixed-deficiency recognizable UNSAT core;
\item a genuinely OR-clause-specific obstruction, surviving the 3XOR-SLICE-EASY hostile control.
\end{enumerate}

This resembles a quiet-planting requirement: the contradiction must be globally detectable by the target function but locally hidden from low-complexity monotone statistics.

\section{Literature implications}

The fixed-deficiency literature is useful primarily as a negative guide: minimal unsatisfiable formulas with fixed clause-variable difference have polynomial recognition algorithms, so they are unsuitable as the entire positive test family under P/POLY-POSITIVE-BARRIER.

Random-CNF proof-complexity work above the satisfiability threshold supplies large-core and expansion structure, but embedding a high-density core into a density-one formula creates detectable degree imbalance unless the embedding is degree-balanced.

The next constructive target should therefore be a degree-balanced large-core transformation rather than simple variable padding.

\section{Next target}

\[
\boxed{\text{DENSITY1-BALANCEDCORE1}}
\]
Construct or rule out a polynomial-size transformation that maps a known large-core unsatisfiable 3-CNF to an exact-\(q\)-clause, \(q\)-variable proper 3-CNF while:
\begin{enumerate}
\item preserving UNSAT;
\item matching the target Poisson-three variable-degree profile to low order;
\item avoiding a fixed-deficiency recognizable image class;
\item preserving enough expansion/large-core structure to support a SAT-specific sunflower argument.
\end{enumerate}

The first hostile control for every construction is the explicit monotone degree-profile separator above.

\section{Status}
No unrestricted lower bound is proved. The note retires one natural positive distribution and adds a general recognizability barrier. \(P\ne NP\) remains OPEN.

\begin{thebibliography}{3}
\bibitem{Valiant1986}
L. G. Valiant, Negation is Powerless for Boolean Slice Functions, \emph{SIAM Journal on Computing} 15(2):531--535, 1986.
\bibitem{Szeider2004}
S. Szeider, Minimal unsatisfiable formulas with bounded clause-variable difference are fixed-parameter tractable, \emph{Journal of Computer and System Sciences} 69(4):656--674, 2004.
\bibitem{FleischnerEtAl2002}
H. Fleischner, O. Kullmann, S. Szeider, Polynomial-time recognition of minimal unsatisfiable formulas with fixed clause-variable difference, \emph{Theoretical Computer Science} 289(1):503--516, 2002.
\end{thebibliography}
\end{document}
